\usetikzlibrary{calc}

% Arrows point from the current cell back to its predecessor.
%
% Namespaced \aln... rather than the generic \drawup/\drawleft/\drawmatch,
% which are plausible names for any document that draws arrows. Older generated
% output defines those three itself with \newcommand; keeping them here made
% \usepackage{alignment-macros} collide with such a file on load.
\newcommand\alnup[3]{%
  \draw[#3] ($(#1,#2) - (0.15,0)$) -- ($(#1,#2) - (0.85,0)$);
}
\newcommand\alnleft[3]{%
  \draw[#3] ($(#1,#2) - (0,0.15)$) -- ($(#1,#2) - (0,0.85)$);
}
\newcommand\alnmatch[3]{%
  \draw[#3] ($(#1,#2) - (0.15,0.15)$) -- ($(#1,#2) - (0.85,0.85)$);
}

% ---------------------------------------------------------------------------
% Visualisation palette — shared with the web application.
% Hue encodes what the move means; it is the same in linear and affine mode:
%   M / diagonal   — substitution (match or mismatch)
%   X / vertical   — gap in the top sequence
%   Y / horizontal — gap in the left sequence
% Weight and saturation encode optimality: dp-* are candidate moves the
% recurrence considered, opt-* are moves on an optimal path.
% Only opt-* carry an arrowhead. A cell's predecessors are always up and/or
% left, so on a candidate arrow the head is redundant with the geometry and
% only invites the question "why does this point backwards?". On the optimal
% path the head marks a traversal the reader actually follows, and reinforces
% the same distinction that weight and saturation already carry.
% ---------------------------------------------------------------------------
% Nothing here may assume a white page. Two knobs carry the document's own
% colours; every other colour is mixed from a hue toward one of them, so the
% picture inverts correctly on a dark theme instead of glowing against it:
%   alnbg — the page. Score patches fill with it; pale candidate arrows mix
%           toward it, so "washed out" always means "closer to the page".
%   alnfg — the ink. Saturated score text mixes toward it.
% Under beamer both follow the theme; elsewhere they default to black on white,
% which reproduces the previous hardcoded output exactly.
% \ifdefined rather than \@ifclassloaded because this file is also inlined
% verbatim into standalone output, where @ is not a letter.
\colorlet{alnbg}{white}
\colorlet{alnfg}{black}
\ifdefined\usebeamercolor
  \AtBeginDocument{%
    \usebeamercolor{background canvas}%
    \colorlet{alnbg}{bg}%
    \usebeamercolor{normal text}%
    \colorlet{alnfg}{fg}%
  }
\fi

% Styles resolve the mix at use time, so a document may \colorlet either knob
% (or override a style outright) at any point and the picture follows.
% Text sizes live here, not in the renderer, so a deck can enlarge the scores
% without regenerating the picture:  \tikzset{aln-score/.append style={scale=1}}
% The three affine state scores share one cell at +-0.3 offsets, so they must
% stay smaller than the single linear score or three-digit values collide.
% NB: TikZ's scale key is multiplicative, so the shared base must not set one
% or a style building on it would compound the two factors.
\tikzset{aln-score-base/.style={fill=alnbg,inner sep=1.5pt}}
\tikzset{aln-score/.style={fill=alnbg,scale=0.7}}
\tikzset{aln-state-score/.style={aln-score-base,scale=0.45}}
\tikzset{aln-state-score-main/.style={aln-score-base,scale=0.5}}
\tikzset{aln-seq-label/.style={scale=1}}
\tikzset{aln-grid/.style={color=alnfg!25!alnbg}}

\definecolor{vizM}{HTML}{56B4E9}
\definecolor{vizX}{HTML}{DC2626}
\definecolor{vizY}{HTML}{009E73}

\tikzset{dp-M/.style={draw=vizM!40!alnbg,thick}}
\tikzset{dp-X/.style={draw=vizX!40!alnbg,thick}}
\tikzset{dp-Y/.style={draw=vizY!40!alnbg,thick}}
\tikzset{opt-M/.style={->,>=latex,draw=vizM,fill=vizM,very thick}}
\tikzset{opt-X/.style={->,>=latex,draw=vizX,fill=vizX,very thick}}
\tikzset{opt-Y/.style={->,>=latex,draw=vizY,fill=vizY,very thick}}

% Dimmed variants, used by the affine decomposition overlays to push the
% two states that are not being discussed into the background.
\tikzset{dp-M-dim/.style={draw=vizM!15!alnbg,thin}}
\tikzset{dp-X-dim/.style={draw=vizX!15!alnbg,thin}}
\tikzset{dp-Y-dim/.style={draw=vizY!15!alnbg,thin}}
\tikzset{opt-M-dim/.style={->,>=latex,draw=vizM!25!alnbg,fill=vizM!25!alnbg,thick}}
\tikzset{opt-X-dim/.style={->,>=latex,draw=vizX!25!alnbg,fill=vizX!25!alnbg,thick}}
\tikzset{opt-Y-dim/.style={->,>=latex,draw=vizY!25!alnbg,fill=vizY!25!alnbg,thick}}

% Affine gap state styles — score text, matching the arrow palette.
\tikzset{text-M/.style={text=vizM!80!alnfg}}
\tikzset{text-X/.style={text=vizX!80!alnfg}}
\tikzset{text-Y/.style={text=vizY!80!alnfg}}
\tikzset{text-M-dim/.style={text=vizM!30!alnbg}}
\tikzset{text-X-dim/.style={text=vizX!30!alnbg}}
\tikzset{text-Y-dim/.style={text=vizY!30!alnbg}}
